\documentclass[a4paper,12pt,twoside]{article} %Dokumtentklasse Artikel, A4-Format, Schriftgr��e 12pt, ACHTUNG: oneside einseitig, twoside f�r doppelbedruckte Seiten

\usepackage[ngerman]{babel} % �bersetzung der sprachspezifischen Befehle im Dokumente (z.B. Inhaltsverzeichnis statt table of content), richtige Trennung nach der deutschen Rechtschreibung, erleichterte Eingabe von Sonderzeichen
\usepackage[ansinew,latin1]{inputenc} % Unterst�tzung erweiterter Eingabe-Zeichens�tze z.B. �, �, � usw.
\usepackage[T1]{fontenc} % f�r Umlaute
\usepackage{eurosym} % Euro-Symbol als \euro einbinden
\usepackage{graphicx} % Paket um Bilder einzuf�gen
\usepackage[paper=a4paper, left=35mm, right=30mm, top=30mm, bottom=20mm]{geometry}  % Satzspiegel f�r wissenschaftl. Arbeit!
\usepackage{amsmath} % Paket f�r mathematische Formeln
%\usepackage{mdwlist} % Zeilenabstand bei Aufz�hlungen verringern
%\usepackage[bottom,hang]{footmisc} %Einr�ckung bei Fu�noten
%\setlength{\footnotemargin}{0pt} %Einr�ckung bei Fu�noten
%\usepackage[hang]{caption2} % mehrzeilige Captions ausrichten
%\usepackage{float} % u.a. genaue Platzierung von Objekten
\usepackage{multirow}
\usepackage{tabularx} %Tabellen sind so gro� wie maximal zur Verf�gung stehende Breite des Texts
%================

\renewcommand{\familydefault}{\sfdefault} % serifenlose Schrift
\renewcommand{\labelitemi}{-} % Spiegelstriche statt Punkten bei Auflistungen
\renewcommand{\contentsname}{Inhalt} % "`Inhaltsverzeichnis"' in "`Inhalt"' umbenennen
%\renewcommand{\bibname}{Literatur} % Literaturverzeichnis umbenennen
\renewcommand{\figurename}{Abb.} % Abbildung umbenennen

\linespread{1,5}\selectfont % aktueller Zeilenabstand wird mit dem Zahlenwert multipliziert

\begin{document}
%=======================================================================
\begin{titlepage}
\vspace*{-2.26cm}

\begin{center}
{ \vspace*{1 cm}
 {\textbf{\Large Praktikumsbericht}\\	
	\vspace*{2 cm}	
 \textbf{Vorlage zum Verfassen einheitlicher Dokumente mit inhaltlichen und stilistischen Hinweisen}
 }
}\end{center}

\vfill
\small{
\begin{tabbing}
\hspace*{6cm}\= \kill 
 Autor:				\> Anna G�nter\\
 Datum der Abgabe: 			\> 26.~Februar~2013\\
 Ort der Abgabe: 				\> Lippstadt\\
 Hochschule:	\> Hochschule Hamm-Lippstadt\\
 Studiengang: \> Mechatronik 7. Semester\\
 Betreuender Professor: \> Prof. Dr.-Ing. Ulrich Schneider\\
\end{tabbing}
}
\end{titlepage}
%===========================================================================
\newpage
%========================================================================
\pagenumbering{Roman} % r�mische Zahlen als Nummerierung von Deckblatt bis Textteil
\noindent\Large\textbf{Sperrvermerk}\par
\vspace*{0.5 cm}
\normalsize\noindent{Diese Arbeit einschlie�lich aller ihrer Teile ist urheberrechtlich gesch�tzt.
Jede Verwertung au�erhalb der engen Grenzen des Urheberrechtsgesetzes ist ohne Zustimmung der Verfasserin unzul�ssig.}\\

\vspace*{1 cm}
\noindent\Large\textbf{Restricting note}\par
\vspace*{0.5 cm}
\normalsize\noindent{This thesis including all its parts is copyrighted.
Any utilisation outside of the strict limits of the copyright law is illegitimate without the permission of the editor.}\\

\vfill
Lippstadt, Februar 2013 \hfill \emph{Anna G�nter} % handschriftliche Unterschrift hierhin!!

%========================================================================
\newpage
\tableofcontents
%========================================================================
\newpage
\noindent
\textbf{\Large{Kurzfassung}}

\vspace*{0.5 cm}

\noindent
Das vorliegende Dokument mit dem Thema "`Vorlage zum Verfassen einheitlicher Dokumente mit inhaltlichen und
stilistischen Hinweisen"' soll als Grundlage zum Erstellen von wissenschaftlichen Arbeiten an der Hochschule Hamm-Lippstadt (im Folgenden als HSHL bezeichnet) dienen.
Sie enthält inhaltliche, stilistische und formale Vorgaben, die unter Verwendung der angegebenen Quellen im Literaturverzeichnis zusammengetragen worden sind.

Das Ziel ist es, den Studierenden der HSHL eine einheitliche Basis für das Verfassen von wissenschaftlichen Arbeiten wie z.\,B. Praktikumsberichten, Projektarbeiten und der Bachelorarbeit zu geben.
Die Vorlage dient sowohl als Hilfe zum Verfassen einer wissenschaftlichen Arbeit an sich als auch zur ihrer Umsetzung mit dem Textsatzprogramm \LaTeX, welches zunehmend in naturwissenschaftlichen Studiengängen und insbesondere an der HSHL verwendet wird.

Zur Veranschaulichung sind im LaTeX-Quellcode dieser wissenschaftlichen Arbeit die geforderte Formatierung und die verwendeten Latex-Befehle mit Kommentaren zu entnehmen\footnote{Die Kurzfassung soll die Relevanz und Zielsetzung der Ausarbeitung herausstreichen und die wichtigsten Inhalte in kurzer und prägnanter Form enthalten (vgl. \cite{Karm2011}, S. 45).}.

%========================================================================
\newpage
\noindent
\textbf{\Large{Abstract}}

\vspace*{0.5 cm}

\noindent
"`An abstract is a succinct summary of the entire paper.
The abstract should briefly describe the question posed in the paper, the methods used to answer the question, the results obtained, and the conclusions.
It should be possible to determine the major points of a paper by reading the abstract"' (\cite{Karm2011}, S. 41).

Immer häufiger wird für wissenschaftliche Ausarbeitungen auch die Verfassung einer englischsprachigen Zusammenfassung, des sogenannten \textsl{Abstracts}, verlangt (vgl. \cite{Karm2011}, S. 47).
Das \textsl{Abstract} beinhaltet den Text der Kurzfassung in englischer Sprache.
Die Zusammenfassung sowie das \textsl{Abstract} sollten sich auf maximal eine Seite belaufen und die wesentlichen Inhalte sowie Ergebnisse herausstellen.
Sie werden nach dem Inhaltsverzeichnis platziert und besitzen keine Kapitelnummerierung.

%==============================================================================
\newpage
%=======================================================================
\chapter*{Vorwort}
%=======================================================================
Das Vorwort beinhaltet pers�nliche Zeilen des Autors zur Arbeit.
Hier wird die Danksagung an die Betreuer und alle, die an der Arbeit behilflich waren, untergebracht (vgl. \cite{Karm2011}, S. 45).

Alle Kapitel vor der Einleitung besitzen keine Nummerierung im Titel.
Die Seitenz�hlung beginnt mit dem Deckblatt, jedoch wird auf dem Deckblatt und dem Sperrvermerk keine Seitenzahl angezeigt.
Erst ab dem Inhaltsverzeichnis werden die Seiten in r�mischen Zahlen fortlaufend nummeriert.
Der Textteil , der mit der Einleitung beginnt, und alle Teile danach werden mit Dezimalziffern beginnend mit der Seite "`1"' fortlaufend durchnummeriert.
Die eidesstaatliche Erkl�rung am Ende der Arbeit enth�lt ebenfalls keine Seitenangabe (vgl. \cite{Balz2008}, S. 185).

Ein Vorwort ist erst beim Verfassen einer Dissertation oder eines Literarischen Werkes �blich.


%========================================================================
\newpage
\noindent
\textbf{\Large{Abbildungs- und Tabellenverzeichnis}} 
\vspace*{0.5 cm}
%
\begin{table*}[h]
	\begin{tabularx}{\textwidth}{lXr}
		%===========================Abbildungen=============================================================================================
		Abbildung \ref{fig:KompassSensor}: & Compass Sensor \dotfill & \pageref{fig:KompassSensor}\\ % Referenz auf Bildnummer und Seite auf der das Bild erscheint
		Abbildung \ref{fig:Diagramm1}: & Beispielhaftes Diagramm zur Verdeutlichung der Anordnung \dotfill & \pageref{fig:Diagramm1}\\ 
		%===========================Tabellen=============================================================================================
	& & \\
		Tabelle \ref{tbl:Laplace}: & Laplace Transformation \dotfill & \pageref{tbl:Laplace}\\ % Referenz auf Tabellennummer und Seite auf der die Tabelle erscheint
		\end{tabularx}
	\label{tab:Abbildungsverzeichnis} 
\end{table*}

%==========================================================================
\newpage
\noindent
\textbf{\Large{Abk�rzungsverzeichnis}} 
\vspace*{0.5 cm}
%
\begin{table*}[h]
	\begin{tabularx}{\textwidth}{l|X}
	Abk�rzung & Bedeutung\\ \hline
	HSHL & Hochschule Hamm-Lippstadt\\
		\end{tabularx}
	\label{tab:Abkuerzungsverzeichnis}
\end{table*}


%=======================================================
\newpage
\pagenumbering{arabic} % arabische Nummerierung der Seiten ab Hauptteil
\setcounter{page}{1} % setzt die Seitennummerierung ab 1 an
%=======================================================================
\chapter{Einleitung}
\label{sec:Einleitung}
%=======================================================================
"`In einem Praktikum erhält der Studierende die Gelegenheit, theoretisches Wissen aus dem Studium unter realen Bedingungen auf den Prüfstand zu bringen.
Die hier gemachten Erfahrungen werden in einem Praktikumsbericht zusammengefasst.
Auch wenn ein Praktikumsbericht im eigentlichen Sinne keine wissenschaftliche Arbeit darstellt, können dennoch die gleichen Anforderungen an Form und Bearbeitung gelten.
[...] Der Aufbau eines Praktikumsberichtes kann dem einer wissenschaftlichen Arbeit entsprechen, wenn er auch nicht als solche zu werten ist.
Die Verwendung von Deckblatt und Inhaltsverzeichnis vor dem eigentlichen Textteil ist auch bei dieser Art von Arbeiten sinnvoll"' (\cite{Balz2008}, S. 157 f.).

Das Deckblatt einer wissenschaftlichen Arbeit ist der erste Eindruck, der visuell wahrgenommen wird und sollte daher ansprechend und übersichtlich gestaltet sein.
Es sollte schlicht und sachlich gehalten werden ohne flächenfüllende Grafiken und Abbildungen zu verwenden (vgl. \cite{Balz2008}, S. 186).
Es ist gemeinhin üblich das Logo der Hochschule oder Fakultät auf dem Deckblatt abzubilden, allerdings wird dringend drauf hingewiesen, dass das Logo der HSHL ohne Freigabe nicht auf dem Deckblatt verwendet werden darf.
Das Deckblatt muss laut \cite[S.\,187]{Balz2008}  folgende Elemente enthalten:
\begin{itemize}
\item Name der Hochschule
\item Bezeichnung des Studiengangs
\item Anzahl der Fachsemester
\item Titel und Untertitel der Arbeit
\item Name des Autors
\item Name des Betreuers
\item Ort und Datum der Abgabe 
\end{itemize}
Konsultieren Sie die Rahmenprüfungsordnung Ihre Studiengangs für ggf. abweichende Forderungen.

Inhaltlich wird in der Einleitung die Aufgabenstellung und die Motivation der behandelten Thematik dargestellt.
Es sollen folgende Fragen beantwortet werden:
\begin{itemize}
\item Warum ist das Thema wichtig?
\item Was sind die wesentlichen Inhalte?
\item Welche Fragen will der Bericht beantworten?
\item Wie gehe ich beim Bearbeiten und Beantworten der Frage(n) vor?
\item Für welchen Zweck sollen die Ergebnisse verwendet werden?
\item Welcher technische Fortschritt ergibt sich durch die Ergebnisse? \cite[S.~27]{Karm2011}
\end{itemize}

%=======================================================================
\section{Motivation}
%=======================================================================
\subsection*{11 Gründe warum Sie \LaTeX verwenden sollten.}

\begin{enumerate}
\bfseries
\item \LaTeX\ produziert portable und konsistente Dokumente. \\ \mdseries
    Mit LaTeX  können Sie PS- und PDF-Dateien produzieren. Wenn Sie diese Dateien versenden, brauchen Sie sich keine Sorgen zu machen, dass die Person am empfangenden Ende diese nicht öffnen kann. 
    Um PDF-Dateien zu lesen, benötigt der Empfänger lediglich einen PDF-Reader wie Adobe, Sumatra, Foxit oder Evince. Diese Programme sind Freeware. Mit dem PDF-Format können Sie auch sicherstellen, dass der Empfänger exakt die Formatierung sieht, welche Sie erstellt haben. Beim Austausch von .docx-Dateien kommt es hier häufig zu Problemen, da zwei Nutzer unterschiedliche Software verwenden. Schreiben Sie beispielsweise ein Dokument mit  KWrite und sehen Sie sich das Durcheinander in OpenOffice an. Hinzu kommt, dass als Abgabeformat in Hochschulen häufig das PDF-Format verlangt wird.
\textbf{    \item \LaTeX\ bietet Ihnen alle Vorteile digitaler Formate.}\\
    Ich benutze fast ausschließlich \verb|pdflatex|, um Dokumente zu erzeugen. Das sind bei mir vorrangig Aufgabenzettel und Praktikumsleitfäden. Da die Endergebnisse PDF-Dateien sind, kann ich diese einfach archivieren. So finde ich leicht zu jedem Jahrgang die zugehörigen Aufgaben wieder. Ich arbeite so auch nahezu papierlos und archiviere alles sortiert in Ordnern in SVN oder Sciebo. So kann man auch schnell ein Dokument mit anderen teilen oder als E-Mail versenden.  Text-Dateien und .docx-Dateien sind auch gut, aber LaTeX noch hat einen für mich entscheidenden Vorteil. Dieser bildet den nächsten Punkt.
\textbf{    \item \LaTeX\ bietet hervorragende Mathematik- und Symbolunterstützung.}\\
    Als Studentin oder Student hat man oftmals mit esoterischen mathematischen und logischen Symbolen, griechischen Buchstaben, und Gott weiß was zu tun. \LaTeX\ hat eine erstaunliche Unterstützung hierfür. Wenn Sie beispielsweise eine Herleitung einer Reihenentwicklung, wie der dargestellten Taylorreihe, machen müssen, wollen Sie Ihre Zeit nicht damit verschwenden die Formeln zu setzen.
	\[\sum\limits_{n = 0}^{\inf}\frac {f^{(n)}(a)}{n!}(X - a)^{n}
\]
\LaTeX\ ermöglicht Ihnen die Formel zu "`schreiben"' und zu kopieren, so dass Sie in der kopierten Formel schnell Umformungen einfügen können.

Es ist viel einfacher die Formel in dieser Form einzugeben:    
\begin{lstlisting}
    \sum\limits_{n = 0}^{\inf}\frac {f^{(n)}(a)}{n!}(X - a)^{n}.
\end{lstlisting} 
Auf diese Weise braucht man selbst beim erstellen von komplexen Formeln die Hände nicht von der Tastatur zu nehmen.
\textbf{\item LaTeX bietet präzise Kontrolle über das Layout des Dokuments.}\\
    Ränder, Zeilenabstand, Abstände zwischen Absätzen, Sie können alle Parameter in \LaTeX\ anpassen. Sie brauchen jedoch nicht mit der Maus Lineale und Tabulatoren hin und her zuschieben, sondern können diese Layoutvorgaben zu Beginn des Dokumentes in die \LaTeX-Datei schreiben. \LaTeX\ gibt Ihnen die Möglichkeit, das Layout bis hin zur Anzahl der Pixel zwischen Kuller und den Inhalt von Aufzählungen zu steuern. Dies ist vergleichbar mit CSS und HTML und basiert auf der Grundidee der Trennung von Inhalt und Layout. Dies wird Ihnen erleichtern sich auf den Inhalt Ihrer wissenschaftlichen Arbeit zu konzentrieren, ohne dass Sie durch Layout-Aufgaben immer wieder abgelenkt werden.
\textbf{\item Sie können Makros in LaTeX verwenden.}\\
    Nun, ich bin sicher, Sie haben aus den Informatik-Vorlesungen eine wichtigen Grundsatz mitgenommen: "`Wiederholungen sind böse und sollten durch Funktionen, Schleifen oder Makros ersetzt werden."'  Verwenden von Makros spart Ihnen Arbeit und erleichtert die Wiederverwendbarkeit bei gleichzeitig Minimierung von Fehlern. \LaTeX\ ermöglicht die Verwendung von Makros in Ihrem Dokumenten. Dieses ist definitiv eine großartige Erweiterung. Das folgende Beispiel zeigt ein solches Makro, dem ein Argument (im Beispiel der Name "`Peter"') übergeben wird:

\begin{lstlisting}
    \newcommand {\helloworld} [1] {
    % Dies ist ein Makro dem ein Argument übergeben wird
    Hallo Welt an {#1}!
    }
    
    % Das Makro wird so aufgerufen
    \helloworld{Peter}
    
    % ...und ergibt nach dem Compilieren
    Hallo Welt an Peter!
\end{lstlisting} 
\textbf{\item Sie können in \LaTeX-Kommentare verwenden.}\\
    Wie Sie im obigen Quelltextauszug vielleicht bemerkt haben, können Sie Kommentare in LaTeX nutzen. Ich halte dies für sehr hilfreich, da ich Kommentare und Notizen für mich oder andere Autoren in die Dokumente einfügen kann. Diese Kommentare sind im PDF-Dokument jedoch nicht sichtbar. Der Leser merkt somit nichts von den internen Notizen. 

\begin{lstlisting}   
    % OpenOME 2.1 Tutorial
		% Noch zu erledigen:
    % - Eclipse Setup-Prozess für Einsteiger
    % - Newsgroup Einbindung
    % - Build-Anleitung
    % - Links zu externen Quellen
    % - Bugfix: setPreferredSize Methode! (siehe bug report # 1211)
\end{lstlisting} 

    Wie Sie sehen können, ist es dies sehr praktisch. Ich verwende häufig Kommentare, um meine Gedanken zu ordnen oder um Textausschnitte auszukommentieren, die ich später überarbeiten und daher nicht löschen möchte.
\textbf{\item Sie können für \LaTeX-Dateien bequem eine Versionskontrolle wie z.\,B. SVN verwenden.}\\
    LaTeX-Dateien (mit der Endung .tex) sind reiner ASCII-Klartext. Docx-Dateien hingegen sind die binäre Dateien. Dies bedeutet, dass Sie die \LaTeX-Dateien mit eine Versionskontrolle einfach verwalten können. Die gängigen SVN-Operationen wie "`revision"', "`log"' "`diff"' und "`merge"' lassen sich somit auf \LaTeX-Dateien anwenden. Ich bin sicher, Sie wissen, warum Sie eine Versionskontrolle verwenden sollten, und werde Sie daher an dieser Stelle nicht mit Details langweilen. 
\textbf{\item \LaTeX\ hat hervorragende Bibliotheken.}\\
    Sie können den Funktionsumfang bequem mit \LaTeX-Pakete und Bibliotheken erweitern. Hier gibt es Vorlagen für Dokumenten, Veröffentlichungen, Notizen, Präsentationsfolien, Grafiken uvm.
\textbf{\item Es gibt eine Menge Unterstützung und Ressourcen für Leute, die LaTeX lernen wollen.}
    Neben Google gibt es auch eine kostenloses und phänomenales \href{http://de.wikibooks.org/wiki/LaTeX-Kompendium}{Lehrbuch} für Anfänger, die LaTeX erlernen möchten \cite{Kompendium2013}. Es gibt auch eine Menge von Diskussionsforen im Internet zum Thema. Hier finden Sie für jede Frage eine passende Antwort.
\textbf{\item \LaTeX\ ist innerhalb der wissenschaftlichen Gemeinschaft bekannt.}\\
    Ernsthaft, wo haben Sie von LaTeX das erste Mal gehört? Im zweiten Semester an der Hochschule/Uni, oder? LaTeX ist recht prominent in der Wissenschaft. Wenn Sie eine Karriere in der Forschung machen möchten, werden sich LaTeX-Kenntnisse auf jeden Fall auszahlen. Ihre Professor/innen und Mitarbeiter/innen sind wahrscheinlich sehr kompetent und könnte Ihnen bei Fragen oder Problemen helfen.
    \item \textbf{Last but not Least ist \LaTeX\ eine plattformunabhängige Freeware und bietet eine perfekte Buchsatzqualität.}
\end{enumerate}
    
Dieser Abschnitt lehnt sich an den Artikel "`10 Reasons Why Computer Science Students Should Learn LaTeX"' an \cite{Hajono2010}. 

%=======================================================================
\section{Aufgabenstellung}
%=======================================================================
In diesem Abschnitt sollte die Aufgabenstellung konkret beschrieben werden. Ihr Ergebnis wird am Erfüllungsgrad der Aufgaben gemessen.

Laborversuche müssen grundsätzlich so dokumentiert werden, dass sie "`reproduzierbar"' sind.
Das bedeutet, dass alle Informationen angegeben werden müssen, um die Versuche unter den exakt gleichen Bedingungen erneut durchführen zu können.
Es muss gewährleistet werden, dass andere Personen bei einer Wiederholung der Untersuchung zu den gleichen Ergebnissen kommen.
Ein hoher Grad an Reliabilität und Messgenauigkeit ist wichtig, d.\,h. dass die Messergebnisse zuverlässig und stabil sind \cite[S.~22~f.]{Balz2008}.
Aus diesem Grund muss die Aufgabenstellen präzise formuliert und die Versuchsvorbereitung, sowie der Aufbau, die Durchführung und Auswertung detailliert geschildert werden.
%=======================================================================
\section{Projektplan}
%=======================================================================
Besonders bei ingenieurwissenschaftlichen Themen wird das Vorgehen detailliert geplant. Stellen Sie Ihren Projektplan als Gantt-Chart (vgl. Bild \ref{fig:Gantt}) vor und heben Sie notwendige Planänderungen hervor. Die Arbeitspakete sollten nachverfolgbar sein, bei Gruppenarbeiten einer Ressource zugeordnet, Vorgänger sowie Nachfolger und Meilensteine aufweisen. Die Zeiteinheit ist sinnvoll zu wählen. Arbeiten Sie an einem Projekt ganztägig oder stundenweise entscheidet die Zeiteinheit.
\begin{figure}[htbp]
	\centering
		\includegraphics[width=0.90\textwidth]{Bilder/Gantt.png}
	\caption{Beispiel für ein Gantt-Chart}
	\label{fig:Gantt}
\end{figure}


%=======================================================================
\section{Gliederung}
%=======================================================================
Die Gliederung verknüpft die Aufgabenstellung mit der Lösung. Ein Beispiel könnte lauten: "`Kapitel \ref{sec:BasicRulez} beschreibt die theoretischen Grundlagen usw."'



%========================================================================
\newpage
\section{Grundlegender Aufbau}
Auf eine �berschrift folgt immer ein Absatz mit Text, \textbf{nie unmittelbar} eine �berschrift niedrigeren Ranges (vgl. \cite{Balz2008}, S. 197).\footnote{Hervorhebungen werden durch Kursivschrift oder halbfette Schrift realisiert. Vermieden werden soll Hervorhebung durch Unterstreichen, gr��ere Schrift oder Gro�buchstaben (vgl. \cite{Balz2008}, S. 182). Ebenfalls ung�nstig ist die farbliche Hervorhebung von W�rtern oder �berschriften.}
Daher muss von der Einleitung eine schl�ssige �berleitung zum Hauptteil und den anstehenden Unterkapiteln erfolgen.
Falls mehrere Versuche innerhalb des Praktikums durchgef�hrt werden, sollte an dieser Stelle kurz darauf eingegangen werden.

Ein Hauptkapitel (Kapitel mit nur einer Z�hlnummer) beginnt jeweils auf einer neuen Seite.
Innerhalb der Gliederung ist zu beachten, dass Gliederungsebene vier nicht �berschritten werden sollte.
Die Unterteilung soll m�glichst in der dritten Stufe enden (vgl. \cite{Hering2009}, S. 12).
Zum Beispiel soll eine Gliederungsebene "`2.2.1.3"' vermieden werden, also wird die Untergliederung bei Ebene "`2.2.1"' beendet.
Desweiteren muss darauf geachtet werden, dass ein Oberpunkt der Gliederung mindestens zwei Unterpunkte besitzt muss (vgl. \cite{Karm2011}, S. 53).
Dennoch sollen die Z�hlnummern innerhalb einer Hierarchiestufe nicht gr��er als neun werden (vgl. \cite{Hering2009}, S. 12).
Die jeweils �bergeordneten Kapitel sollen einen �berblick �ber die Inhalte geben, die in den Unterkapiteln thematisiert sind.
In den �berblicken sollte jedoch kein zus�tzliches Wissen vermittelt werden (vgl. \cite{Balz2008}, S. 196).

In diesem Dokument wird ein beispielhaftes Inhaltsverzeichnis f�r einen Praktikumsbericht aufgezeigt, der u.a. im Praktikum Grundlagen der Elektrotechnik an der HSHL angefertigt werden muss.
Die Inhalts�bersicht und das Inhaltsverzeichnis sollen einen �berblick �ber die Systematik der Arbeit geben und einen "`ersten Eindruck"' vermitteln.
Man sollte sich bewusst sein, dass das Inhaltsverzeichnis die erste Information ist, die ein Leser von der Arbeit aufnimmt (vgl. \cite{Karm2011}. S. 49).
Die Gliederung soll bereits als "`roter Faden"' dienen.
F�r den Leser ist die Gliederung der wichtigste Schl�ssel zum Erfassen des Inhalts.
Wissenschaftliche Arbeiten erfordern ein hohes Ma� an Ordnung und Logik, welche sich auch in der Gliederung widerspiegeln muss (vgl. \cite{Hering2009}, S. 9 ff.).

\newpage
\noindent
�blicherweise gliedert sich der inhaltliche Aufbau in Einleitung, Hauptteil und Schluss.
Die Einleitung beinhaltet die Aufgabenstellung und die Motivation der wissenschaftlichen Arbeit (s. Einleitung, S. \pageref{sec:Einleitung}).
Im Hauptteil werden die theoretischen Grundlagen erl�utert und der L�sungswegs sowie die Ergebnisse dargestellt.
Die Ergebnisse werden im Schlussteil in Kurzform wiedergegeben und kritisch diskutiert.
Weitere Untergliederungen dieser drei Hauptpunkte k�nnen je nach Art der wissenschaftlichen Arbeit und den entsprechenden Anforderungen umgesetzt werden.
In dieser wissenschaftlichen Arbeit werden die ben�tigten Gliederungspunkte f�r einen Praktikumsbericht aufgezeigt.

Ein Praktikumsbericht, der die Beschreibung einen Versuchs beinhaltet, sollte unbedingt folgende Unterpunkte beinhalten:
\begin{enumerate}
\item Theoretische Grundlagen 
\item Aufgabenstellung mit Versuchsbeschreibung
\item Versuchsvorbereitung (ben�tigter Materialien und Beschreibung des Versuchsaufbaus)
\item Versuchsdurchf�hrung
\item Versuchsauswertung (vgl. \cite{Hering2009}, S. 18)
\end{enumerate}

Als allgemeine Gliederung f�r ingenieurwissenschaftliche Arbeit (z.\,B Bachelor- oder Masterarbeit) empfiehlt sich die Nachfolgende. 
\begin{enumerate}
	\item Titelseite
	\item Inhaltsverzeichnis
	\item Einleitung (Einf�hrung in die Aufgabenstellung)
  \item Theoretische Grundlagen
	\item L�sungsweg ggf. L�sungsalternative
  \item Ergebnis
  \item Zusammenfassung (mit anschl. Diskussion der Ergebnisse) und Ausblick
\end{enumerate}
Quellenverzeichnis
Anhang
Eidesstattliche Erkl�rung 


%========================================================================
\newpage
%=======================================================================
% Titelseite
%=======================================================================
\chapter{Theoretische Grundlagen}
\label{sec:BasicRulez}

\begin{figure}[h]
	\centering
		\includegraphics[width=0.80\textwidth]{Bilder/Python_God.jpg}
	\caption{Es gibt Regeln an die man sich halten sollte \cite{Python83}}
	\label{fig:Python_God}
\end{figure}
In diesem Unterkapitel "`Theoretische Grundlagen"' soll ein Überblick über die unterschiedlichen Definitionen und Ansätze zum jeweiligen Thema und eine Abgrenzung zu verwandten Themen und Begriffen erfolgen \cite[S.~27]{Karm2011}.
Wissenschaftliche Arbeiten sollen Klarheit beim Leser hervorrufen.
Das bedeutet, dass der Leser die beschriebenen Sachverhalte exakt in der beabsichtigten Weise ohne Rückfragen verstehen muss.
Welche Definitionen sind relevant um den Bericht ohne weitere Hintergrundkenntnisse in dem Themenbereich zu verstehen?
Inhalte müssen logisch gegliedert und klar nachvollziehbar sein.
Es ist ein hohes Maß an Systematik, Ordnung, Logik und Klarheit erforderlich \cite[S.~1~ff.]{Hering2009}.

"`Zwischendurch müssen Sie immer wieder einmal die gedankliche Position des Lesers einnehmen und sich fragen:
\begin{itemize}
	\item Wann braucht der Leser welche Information?
	\item Fällt das hier verwendete Bild 'vom Himmel'?
	\item Soll ich mich auf die Gliederung beziehen, eine Zwischenzusammenfassung machen oder den neuen Dokumentteil aus übergeordneter Sicht ankündigen?
	\item Ist die Aufteilung der Informationen logisch und nachvollziehbar?"' \cite[S.~29]{Hering2009}.
\end{itemize}

Die nachfolgenden Grundregeln sind strikt zu beachten. 

% Zur besseren Übersicht können Abschnitte in Dateien ausgelagert werden.
\input{Kapitel/Aufbau}
\input{Kapitel/Stilistische_Anforderungen}
\input{Kapitel/Zitieren}
\input{Kapitel/Bilder}





%========================================================================
\newpage
\subsection{Aufgabenstellung}


Im nachfolgenden Kapitel werden einige stilistische und inhaltliche Anforderungen an den Text einer wissenschaftlichen Arbeit aufgezeigt.

%========================================================================
\subsection{Stilistische Anforderungen}
Die "`Ich"'- oder "`Wir"'-Perspektive soll im gesamten Textteil vermieden werden.
Stattdessen sollte aus der dritten, neutralen Person oder in Passivkonstruktionen geschrieben werden (vgl. \cite{Hering2009}, S.130).
Die wissenschaftliche Arbeit darf nicht im Sinne eines Erlebnisberichtes sondern als sachorientierte Dokumentation aufgefasst werden. 
Eine Ausnahme besteht im Schlussteil, in dem der Autor eine kritische Auseinandersetzung mit der Thematik einbringen darf.

Ein wesentliches Kriterium f�r die Verst�ndlichkeit der Inhalte ist ein wissenschaftlich-technischer, konkreter, klarer und pr�ziser Sprachstil.
Es sollte auf keinen Umst�nden Umgangssprache verwendet werden. 
Folgende Punkte sorgen f�r bessere Textverst�ndlichkeit:
\begin{itemize}
	\item Kurze S�tze bilden.
	\item M�glichst nur einen Hauptsatz mit einem oder zwei Nebens�tzen bzw. zwei mit Semikolon verkn�pfte Haupts�tze verwenden. 
	\item Nicht zu viele Nebens�tze in Klammern sondern als alleinstehende S�tze oder in Fu�noten unterbringen.
	\item "`Welcher, welche, welches"' als Relativpronomen sind veraltet und sollen ersetzt werden durch "`der, die, das"'. 
	\item Ein- und �berleitende S�tze verwenden, um Ihre Leser mit Worten zu f�hren.
Damit k�nnen Sie:\\
\begin{itemize}
	\item die Gliederung bzw. das Inhaltsverzeichnis aufgreifen,
	\item den bisher beschriebenen Sachverhalt zusammenfassen,
	\item zum n�chsten Dokumentteil �berleiten oder
	\item einen neuen Dokumentteil einleiten.
\end{itemize}
	\item Anschaulich formulieren, in Bildern sprechen!
	Analogien, Metaphern und Vergleiche erzeugen Assoziationen beim Leser.
	Dadurch kann der Leser �hnlichkeiten und Abweichungen von dem ihm bereits bekannten Wissen leichter erkennen (\cite{Hering2009}, S. 129).
	\item Konjunktive vermeiden.
	\item Apostrophe sparsam im Text verwenden.
	\item Indirekte statt direkte Zitate verwenden.
	\item Wiederholungen von W�rtern in S�tzen vermeiden.
	\item F�llw�rter, d.h. W�rter die beim Weglassen keinen Einfluss auf den Inhalt haben, vermeiden.
	\item Inhaltliche Bl�cke im Text bilden und diese durch Einr�ckungen voneinander trennen.
\end{itemize}

\noindent
Die Erz�hlzeit ist das Pr�sens (Gegenwart).
Das Imperfekt (Vergangenheit) wird nur bei Beschreibung eines fr�her verwendetes Bauteil, Messverfahren usw., eingesetzt.
Es sollten keine emotionsbeladenen und umgangssprachlichen Formulierungen verwendet werden (vgl. \cite{Hering2009}, S. 131).
Zur verst�ndliche Formulierung geh�ren einleitende S�tze am Anfang eines Dokumentteils.
Zusammenfassende und �berleitende S�tze sollten am Ende eines Dokumentteils eingebracht werden.
Diese S�tze k�nnen auch Einsatz zur Ein- und �berleitung bei Bildern und Tabellen finden.
Sie beschreiben z.B. welche Informationen und warum als n�chstes folgen, wie die folgenden Informationen zu bewerten sind, welchen Stellenwert die zuvor beschriebenen Informationen haben usw. (vgl. \cite{Hering2009}, S. 136 f.).

Technische Ger�te sollen grunds�tzlich nach ihrer Funktion benannt werden.
Zum Beispiel ist die Bezeichnung Antriebsrad besser als Zahnrad oder Halteplatte eindeutiger als Platte.
Technische Verfahren werden zur Veranschaulichung funktionsgerecht benannt.
So sollte die Bezeichnung "`Variante 1"' in z.\,B. "`elektrisch-mechanische L�sung"' umbenannt werden (vgl. \cite{Hering2009}, S. 131).
Dabei muss aber beachtet werden, dass keine Worte erfunden werden d�rfen.
Es sollen nur Worte verwendet werden, die im Duden stehen.
Bezeichnungen wie z.B. "`bestromen"' oder "`anflanschen"' sollten in einer wissenschaftlichen Arbeit durch gebr�uchliche Ausdr�cke ersetzt werden.

Bei der Benutzung von Fachbegriffen, Fremdw�rten und Abk�rzungen gilt es eine angemessene Verwendung einzuhalten.
Fremdsprachige W�rter sollten vermieden werden.
Nach M�glichkeit sollen die sinngem��en Bezeichnungen in der deutschen Sprache verwendet werden.
Ausnahme bilden fremdsprachige Fachausdr�cke, bei denen eine �bersetzung ins Deutsche �bertrieben erscheinen w�rde (sogenannte Anglizismen vermeiden).
Bei der Verwendung von Fachbegriffen m�ssen diese bei erstmaliger Verwendung erl�utert werden.
Dies kann sowohl innerhalb des Textes als auch in Fu�noten erfolgen.\footnote{In den Fu�noten k�nnen Anmerkungen untergebracht werden ohne den Lesefluss zu st�ren.
In die Fu�noten geh�ren ausschlie�lich zus�tzliche Informationen und Anmerkungen, die als nennenswert erachtet werden, aber nicht f�r die Argumentation erforderlich sind. Jede Fu�note beginnt mit einem Gro�buchstaben und endet mit einem Punkt (\cite{Balz2008}, S. 197 f.).}
Unbekannte Abk�rzungen werden ebenfalls beim ersten Auftreten definiert (vgl. \cite{Balz2008}, S. 136).
Sobald eine Abk�rzung einmalig eingef�hrt wurde, muss sie konsequent verwendet werden.
Beim Einf�hren einer Abk�rzung wird zuerst der vollst�ndige Begriff genannt und in Klammern auf die Abk�rzung verwiesen wird (s. S. 1 Einf�hrung der Abk�rzung HSHL).
Bei mehr als f�nf verwendeten Abk�rzungen sollte ein Abk�rzungsverzeichnis erstellt werden.

%===================================================================================================
\subsubsection{Materialien}
Die, f�r einen Versuch notwendigen, Materialen werden hier pr�zise mit Artikelbezeichung und Anzahl in Form einer Auflistung aufgef�hrt. 
Folgende Angaben m�ssen unbedingt gemacht werden:

\begin{itemize}
	\item Pr�fmaschine / Versuchseinrichtung mit Hersteller, Typbezeichnung, Inventarnummer usw.
	\item alle jeweils an der Maschine bzw. Einrichtung eingestellten Parameter
	\item alle Messger�te, jeweils mit Hersteller, Typbezeichnung, Inventarnummer, eingestellten Parametern usw.
	\item verwendete Proben mit allen erforderlichen Angaben nach der jeweiligen DIN-, EN-, ISO-Norm bzw. andere Normen je nach Notwendigkeit
	\item bei nicht genormten Versuchen sinngem��e Angaben zur Probenform, zu Versuchsparametern, zu Temperaturen, physikalischen und chemischen Eigenschaften usw.
	\item s�mtliche gemessenen Werte bzw. Versuchsergebnisse mit allen Parametern
	\item angewendete Auswertungsformeln mit vollst�ndigen Quellenangaben usw. (\cite{Hering2009}, S. 24)
\end{itemize}

%================================================================================================
\subsubsection{Der Versuchsaufbau...}
"`...wird in diesem Kapitel beschrieben."'
Dieses Beispiel soll zeigen, dass der Dokumentteil-Titel sprachlich nicht in den darauffolgenden Text einbezogen werden darf.
Der erste Satz des nachfolgenden Textes muss deshalb ein vollst�ndiger Satz sein, der jedoch die Sachaussage des Dokumentteil-Titels aufgreifen bzw. wiederholen darf (vgl. \cite{Hering2009}, S. 14).
Hinter der Dokumentteil-Nummer und dem Dokumentteil-Titel darf nie ein Satzzeichen wie Punkt, Doppelpunkt, Fragezeichen, Ausrufungszeichen usw. stehen (vgl. \cite{Hering2009}, S. 15).


%========================================================================
\include{Versuchsdurchf�hrung}

%========================================================================
\include{Versuchsauswertung}

%========================================================================
\newpage
\section{Zusammenfassung und Ausblick}
"`Ziel des Praktikums ist es, die in einem Studium erworbenen theoretischen Kenntnisse im praktischen Umfeld anzuwenden und individuelle F�higkeiten im Umgang mit diesen zu erwerben. Der Praktikumsbericht enth�lt daher eine ausf�hrliche Reflexion der im Praktikum durchgef�hrten T�tigkeiten und setzt diese in Bezug zu den theoretischen Kenntnissen, welche im Studium erlangt wurden"' (\cite{Balz2008}, S. 158).

Der Schlussteil besteht aus einer kurzen Zusammenfassung und kann auch pers�nliche Bemerkungen, Bewertungen etc. enthalten. Zuerst wird somit das Vorher geschriebene knapp zusammengefasst. Anschlie�end werde die Ergebnisse reflektiert. Die Zusammenfassung setzt sich kritisch mit der Aufgabenstellung auseinander und beantwortet die Fragen:
\begin{itemize}
	\item Was sollte getan werden und was wurde tats�chlich erreicht?
	\item Wo gab es besondere Schwierigkeiten?
	\item Welche Teile der Aufgabenstellung konnten eventuell gar nicht bearbeitet werden und warum?
\end{itemize}

Der Ausblick kann Empfehlungen f�r eine weitere sinnvolle Fortsetzung des Projektes bzw. der Forschungsarbeit enthalten.
Solche Empfehlungen beruhen �blicherweise auf Erkenntnissen, die w�hrend der Bearbeitung des laufenden Projektes entstanden sind (vgl. \cite{Hering2009}, S. 44).\footnote{Es wird darauf hingewiesen, dass die wissenschaftliche Arbeit am Ende unbedingt von einer dritten Person korrekturgelesen werden sollte.
Dadurch k�nnen Fehlerarten wie falsche Wortwahl, Rechtschreibfehler. Zeichensetzungsfehler, Grammatikfehler, Stilfehler, inhaltliche Wiederholungen, Fehler im logischen Gedankenfluss, Layout-Fehler sowie sonstige Fehler aufgedeckt werden (vgl. \cite{Hering2009}, S. 174).
Zudem wird der Autor beim h�ufigen Lesen der eigenen Texte mit der Zeit "`betriebsblind"' gegen�ber den eigenen Formulierungen.
Auch aus diesem Grund sollte von einem Au�enstehenden �berpr�ft werden, ob der Bericht f�r Projektfremde verst�ndlich ist (vgl. \cite{Hering2009}, S. 130).}

%=========================================================================
\newpage
\begin{thebibliography}{}
\label{sec:Literatur}
\bibitem[1]{Karm2011} Karmasin, Matthias/Ribing, Rainer: Die Gestaltung wissenschaftlicher Arbeiten/Ein Leitfaden f�r Seminararbeiten, Bachelor-, Master- und Magisterarbeiten sowie Dissertationen, Wien, 2011, 6. Auflage
\bibitem[2]{Balz2008} Balzert, Helmut/Sch�fer, Christian/Schr�der, Marion: Wissenschaftliches Arbeiten/Wissenschaft, Quellen, Artefakte, Organisation, Pr�sentation, Herdecke/Witten, 2008, 2. Auflage
\bibitem[3]{Kipp2012} HSHL, Skript Wissenschaftliches Arbeiten, Referent: M.A. MLIS Guido Kippelt, Stand 10.05.2012
\bibitem[4]{Hering2009} Hering, Lutz/Hering, Heike: Technische Berichte, Verst�ndlich gliedern, gut gestalten,
�berzeugend vortragen, Vieweg+Teubner | GWV Fachverlage GmbH, Wiesbaden 2009, 6., aktualisierte und erweiterte Auflage
\bibitem[5]{Mess2012} HSHL, Formelsammlung Mess- und Regelungstechnik, WS 2011/12, Prof. Dr.-Ing. Ulrich Schneider
\end{thebibliography}

% Fu�noten bei Zitaten verzichten!
% Quellen in der Reihenfolge, wie sie im Text erscheinen!
% Im Text in normaler Gr��e Nummer in eckiger Klammer einf�gen [1]
% Im Literaturverzeichnis steht [1] und die Erkl�rung
% Literaturverzeichnis MUSS im Inhaltsverzeichnis erscheinen!
  
%==============================================================================================================
\newpage
\appendix
\section{Anhang}
\subsection{Versuchs- und Messprotokolle}\footnote{Im Anhang werden erg�nzende Dokumente untergebracht, die im Textteil wegen ihres Umfangs keinen Platz gefunden haben.
Der Anhang darf auf keinen Fall gr��er als der Umfang des Textteils sein und wird nicht zum Seitenumfang des Textteils hinzugerechnet (vgl. \cite{Balz2008}, S. 198 f.).
Falls die beigelegten Dokumente nicht gebunden werden sollen, k�nnen sie in eine separate Mappe eingelegt oder auf einem Datentr�ger abgegeben werden.
In diesem Fall sollte ein Vermerk wie "`(in separatem Ordner)"' oder "`(als CD-Rom angef�gt)"' im Anhang hinzugef�gt werden.
}
\label{lab:Messprot}

%=========================================================================
\newpage
\pagenumbering{gobble} %Seitennummerierung wird unterdr�ckt
\noindent
Ich versichere hiermit, dass ich die vorliegende Arbeit selbstst�ndig und nur unter Benutzung der angegebenen Literatur und Hilfsmittel angefertigt habe.
W�rtlich �bernommene S�tze und Satzteile sind als Zitate belegt, andere Anlehnungen hinsichtlich Aussage und Umfang unter Quellenangabe kenntlich gemacht.
Die Arbeit hat in gleicher oder �hnlicher Form noch keiner Pr�fungsbeh�rde vorgelegen und ist auch noch nicht ver�ffentlicht.

%An das Ende jeder Bachelorarbeit ist die Eidesstaatliche Versicherung auf eine gesonderte Seite zu setzen -> HSHL-Vorgabe!!

\vfill
Lippstadt, Februar 2013 \hfill \emph{Anna G�nter} % handschriftliche Unterschrift hierhin!!
%=========================================================================
\end{document}